\documentclass[UTF8,zihao=-4]{ctexart}
\usepackage[a4paper,margin=2.2cm]{geometry}
\usepackage{xcolor}
\usepackage{hyperref}
\usepackage{enumitem}
\usepackage{longtable}
\usepackage{array}
\hypersetup{colorlinks=true,linkcolor=teal,urlcolor=teal}
\setlist{itemsep=0.25em,topsep=0.35em}
\title{\bfseries 点子发芽日报 2026-05-27}
\author{点子发芽}
\date{2026-05-27}
\begin{document}
\maketitle
\section*{今日总结}
今天的输入指向一个很具体的提醒：不要只按赛事外壳判断一个场域的价值。中美创客大赛表面上是高校竞赛和青年交流项目，但真正值得观察的是其中的人、项目原型、议题选择和跨文化协作方式，它可能比“比赛”这个词本身更接近一个早期项目筛选入口。
\section*{今日输入}
\begin{longtable}{|p{0.22\linewidth}|p{0.58\linewidth}|p{0.10\linewidth}|}
\hline
\textbf{关键词} & \textbf{补充信息} & \textbf{权重} \\
\hline
中美创客大赛 & 本来以为没啥意思，后来发现里面有些不错的人和项目。 & 3 \\
\hline
\end{longtable}
\section*{今日新知}
\subsection*{中美创客大赛}
\paragraph{简介} 中美创客大赛通常指“中美青年创客大赛”，是围绕青年创新、创客教育和中美人文交流展开的竞赛型项目。它以团队参赛、问题定义、原型设计、现场路演和赛区选拔为主要形式，鼓励青年围绕社区韧性、低碳环保、公共卫生、健康福祉、清洁能源等公共议题提出可落地的产品或服务方案。和纯商业创业比赛相比，它更强调教育属性、跨学科协作、动手实践和国际交流，参赛成果往往处在早期原型阶段，价值不只在获奖，也在于发现有行动力的学生、导师、组织者和初始项目。
\paragraph{最近有什么相关新闻} 可核实的近期信息主要来自 2026 年校内赛和报名通知：部分高校已启动 2026 年中美青年创客大赛报名与校内选拔，公开安排显示分赛区选拔预计持续到 2026 年 6 月下旬，决赛阶段计划在 2026 年 7 月中下旬前后进行，具体时间仍以组委会最终通知为准。
\paragraph{与我相关} 你的补充信息里，关键变化是从“竞赛可能没意思”转为“里面有些不错的人和项目”。这说明它对你不是一个需要泛泛了解的赛事名，而是一个可观察的人才与项目入口：值得看哪些人能把公共议题转成原型，哪些项目虽然稚嫩但有真实需求，哪些团队具备持续迭代的可能。
\paragraph{最小下一步} 整理 3 个你看到的参赛人或项目，用“人、项目、为什么值得继续看”三列各写一句。
\section*{与旧知识的链接}
\begin{itemize}[leftmargin=2em]
\item 连接到“新知孵化工作流”：这类竞赛观察适合被拆成轻量节点，而不是一次性做成完整研究。
\item 连接到“本地 Markdown 知识库”：可以把值得继续看的团队或项目先沉淀为 candidate 节点，后续再决定是否升权。
\item 连接到“点子种子：每日小推进”：明天只需要补一张三列表，而不是展开成大规模调研。
\end{itemize}
\section*{今日发芽点子}
\begin{itemize}[leftmargin=2em]
\item 点子种子：赛事场域观察卡。以后遇到看似普通的比赛、展会或训练营，固定记录“好人、好项目、好问题”三类线索，避免被活动包装影响判断。
\item 点子种子：早期项目雷达。把中美创客大赛这类场域当作早期项目雷达，只保留少量可跟踪对象，后续观察它们是否继续迭代。
\end{itemize}
\section*{参考搜索内容}
\begin{itemize}[leftmargin=2em]
\item \href{https://chinaus-maker.cscse.edu.cn/}{[S1] 官方/机构资料：中美青年创客大赛}
\item \href{https://www.si.sjtu.edu.cn/content/news-detail?params=ba608706abad4c0fa4bf8bdf712a79c1&menuIds=3,26,28}{[S2] 官方/机构资料：2026中美青年创客大赛校内赛启动报名 - si.sjtu.edu.cn}
\item \href{https://yscm.yangtzeu.edu.cn/info/1087/11710.htm}{[S3] 官方/机构资料：关于"2026共创未来——中美青年创客大赛"的报名通知-长江大学传媒学院}
\item \href{http://www.zhjingsai.com/?m=home&c=View&a=index&aid=739}{[S4] 项目主页/机构主页：2025中美青年创客大赛-全国青年大学生竞赛官网-高质量竞赛平台}
\item \href{https://cp.moocollege.com/home/competition/details?competitionId=105634}{[S5] 普通线索：高校学生竞赛网—竞赛质量提升,创新人才培养}
\end{itemize}
\section*{随心记复盘}
\paragraph{主题}
\begin{itemize}[leftmargin=2em]
\item CLI 工具的使用手感
\item Codex 稳定性与配置问题
\item 对模型表现波动的观察
\end{itemize}
\paragraph{讨论} 今天的随心记更像是在比较不同 AI 工具进入日常工作流后的真实摩擦：一边是 Claude CLI 带来的顺手感和期待，一边是 Codex 通过配置后恢复稳定但仍需要观察。这里的张力不是简单换工具，而是在寻找更可靠、更少打断、更贴近自己工作节奏的操作界面。
\paragraph{评价} 这更像一个值得持续观察的长期信号：你已经不只是在试用模型能力，而是在评估工具链是否能成为稳定的个人基础设施。关于“降智”的讨论可以先作为待观察线索，不急着下结论。
\paragraph{留给明天的问题} 明天如果只比较一个维度，你更想验证 CLI 的顺手感，还是验证 Codex 当前配置后的稳定性？
\end{document}
